% Appendix D source fragment. Do not input this file into the six-page paper.
% The overall portfolio wrapper should provide the Appendix D title page and
% appendix numbering before inputting this content.
\section{Design and Implementation Details}
\label{app:design-implementation}

\subsection{Purpose and Design Progression}
This appendix records the implementation detail omitted from the six-page
paper. The project progressed through four modelling questions:

\begin{enumerate}
    \item Can simple point-wise models establish a credible reference?
    \item Does generic nonlinearity improve the selected insertion-loss paths?
    \item Does predicting complete broadband curves jointly provide a better
          representation of frequency structure?
    \item Can one model predict the reciprocal complete complex 12-port
          scattering matrix over the fixed frequency grid?
\end{enumerate}

The sequence is an exploration path, not a claim that each later model is more
accurate. The notation $IL_{7,1}$ denotes insertion loss from input port 1 to
output port 7. Table~\ref{tab:appendix-d-interfaces} shows that the later stages
also change the output scope and data interface. Consequently, comparisons are
used to reveal trade-offs rather than to attribute every difference to one
isolated architectural choice.

\begin{table*}[t]
    \caption{Implemented model interfaces. The batch dimension is omitted.}
    \label{tab:appendix-d-interfaces}
    \centering
    \scriptsize
    \begin{tabular}{@{}llll@{}}
        \toprule
        Stage & Model input & Model output & Main question \\
        \midrule
        Point-wise scalar
          & 10 design values + frequency
          & one $IL_{7,1}$ value
          & basic linear reference \\
        Point-wise vector
          & 10 design values + frequency
          & six insertion-loss values
          & output packaging and nonlinearity \\
        Whole curve
          & 10 design values + fixed frequency features
          & $200\times6$ insertion-loss tensor
          & joint broadband representation \\
        Full matrix
          & 10 design values; internal fixed frequency features
          & $200\times156$ unique real-valued channels
          & complete reciprocal complex response \\
        \bottomrule
    \end{tabular}
\end{table*}

\subsection{Shared Data and Leakage Controls}
All stages used the same 7,030 aligned designs and the project seed 128. The
design-level split contained 4,218 training, 1,406 validation and 1,406 test
designs. Every frequency belonging to one design stayed in the same split.
Training data fitted model parameters and all scaling or feature transforms;
validation data selected configured alternatives and neural checkpoints. Test
designs were excluded from fitting, although their results were inspected
during development, so the final evaluation is retrospective rather than an
untouched confirmation.

The point-wise table contains one row per design and frequency. Its 11 inputs
are ten physical design parameters plus frequency in gigahertz (GHz). The curve
and matrix models retain one sample per design and generate all 200 frequencies
together. This change avoids treating each frequency row as a separate training
sample, but it also means that a point-wise/curve comparison changes more than
the network architecture.

Each model owns its preprocessing so that saved predictions can be reproduced:

\begin{itemize}
    \item Ridge and point-wise neural models standardise their inputs using
          training data. Their insertion-loss targets are also standardised for
          neural training.
    \item Polynomial models expand each input independently and then fit a
          second training-only standardisation. No cross-feature interaction
          terms are generated by this powers-only expansion.
    \item Random Forest uses the original unscaled inputs because tree splits
          depend on ordering rather than feature magnitude.
    \item The whole-curve model standardises the ten design inputs and each of
          its six target channels using pooled training-curve values.
    \item The full-matrix model standardises the design inputs. Each complex
          matrix entry uses one training-set root-mean-square scale shared by
          its real and imaginary parts, preserving their relative scale.
\end{itemize}

\subsection{Non-Neural Point-Wise Models}
\subsubsection{Ridge Regression}
Ridge regression fits a linear mapping while discouraging very large
coefficients. For input matrix $\mathbf X$, target matrix $\mathbf Y$ and
coefficient matrix $\mathbf W$, the fitted objective is
\begin{equation}
    \min_{\mathbf W}
    \left\|\mathbf Y-\mathbf X\mathbf W\right\|_F^2
    +\alpha\left\|\mathbf W\right\|_F^2.
    \label{eq:appendix-d-ridge}
\end{equation}
The scalar version predicts only $IL_{7,1}$. The vector version predicts six
corresponding left-to-right insertion-loss paths. Multi-output Ridge fits one
coefficient vector per target; it does not create a learned coupling between
the outputs. The common selected regularisation value was $\alpha=0.01$.

The scalar model established whether design and frequency features could beat
a pooled training mean. The vector model then established a common six-output
interface for the nonlinear point-wise models. Their identical $IL_{7,1}$
solution is therefore expected and is not evidence that adding outputs improves
the scalar target.

\subsubsection{Polynomial Ridge}
Polynomial Ridge tests explicit smooth nonlinearity without changing the
linear estimator. Each standardised input $x_j$ is expanded independently as
\begin{equation}
    x_j\mapsto[x_j,x_j^2,\ldots,x_j^D].
    \label{eq:appendix-d-polynomial-map}
\end{equation}
Validation selected degree $D=3$ and $\alpha=50$. The expansion contains no
terms such as $x_jx_k$, which keeps the model compact but prevents it from
representing interactions explicitly. This was an intentional intermediate
step between linear Ridge and more flexible learned or tree-based mappings.

\subsubsection{Random Forest}
Random Forest averages 128 regression trees. Each tree recursively divides the
joint design--frequency space into regions and predicts the mean target within
each terminal region. Bootstrap samples diversify the trees before averaging;
the selected implementation makes every input feature available at each split.
Validation retained a minimum of two samples per leaf and unrestricted maximum
depth. The model can represent local interactions without polynomial feature
engineering, but its piecewise predictions do not encode smoothness along
frequency.

\begin{table*}[t]
    \caption{Selected non-neural configurations and their design roles.}
    \label{tab:appendix-d-non-neural}
    \centering
    \scriptsize
    \begin{tabular}{@{}lllll@{}}
        \toprule
        Model & Outputs & Transform & Selected setting & Design role \\
        \midrule
        Scalar Ridge
          & 1 & standardisation & $\alpha=0.01$
          & weak linear baseline \\
        Vector Ridge
          & 6 & standardisation & $\alpha=0.01$
          & common six-output baseline \\
        Polynomial Ridge
          & 6 & powers + standardisation & degree 3, $\alpha=50$
          & explicit smooth nonlinearity \\
        Random Forest
          & 6 & none & 128 trees, leaf $\geq2$
          & local nonlinear interactions \\
        \bottomrule
    \end{tabular}
\end{table*}

These scikit-learn estimators do not have Keras computation graphs. Equations,
the transformation description and Table~\ref{tab:appendix-d-non-neural} are
therefore more informative than forcing them into neural-network-style layer
diagrams.

\subsection{Point-Wise Neural Models}
A multilayer perceptron (MLP) is a feed-forward neural network in which every
unit in one layer connects to every unit in the next. The raw-input MLP accepts
the 11 standardised design--frequency features. Three hidden layers contain
128, 128 and 64 units with rectified linear unit (ReLU) activation, followed by
a linear six-unit output. It has 26,694 trainable parameters.

The Polynomial Neural MLP uses the same hidden and output layers after a fixed
fifth-degree powers-only expansion. Eleven original inputs therefore become 55
expanded inputs before the training-fitted standardisation, increasing the
network to 32,326 parameters. Comparing the two MLPs tests the value of explicit
univariate powers while keeping the learned hidden architecture fixed.

\begin{figure*}[p]
    \centering
    \begin{minipage}[t]{0.46\textwidth}
        \centering
        \includegraphics[width=\linewidth,height=0.72\textheight,
                         keepaspectratio]
          {appendix_d/neural_mlp_topology.png}
    \end{minipage}\hfill
    \begin{minipage}[t]{0.46\textwidth}
        \centering
        \includegraphics[width=\linewidth,height=0.72\textheight,
                         keepaspectratio]
          {appendix_d/polynomial_neural_mlp_topology.png}
    \end{minipage}
    \caption{Saved point-wise MLP topologies. The raw-input model is shown on
    the left and the polynomial-input model on the right. The graphs are
    generated from the selected Keras artifacts; preprocessing occurs outside
    the displayed networks. ``None'' denotes the variable batch dimension.}
    \label{fig:appendix-d-pointwise-mlp}
\end{figure*}

Both networks minimise mean squared error (MSE) in standardised target space
using the Adam optimiser. Training uses batches of 512, an initial learning rate
of $3\times10^{-5}$ and gradient-norm clipping at 0.5. A validation-loss plateau
reduces the learning rate; early stopping restores the best validation
checkpoint. The maximum of 100 epochs is a ceiling rather than a requirement to
retain the last epoch.

\subsection{Whole-Curve Neural Model}
The whole-curve model accepts one ten-parameter design and returns six complete
200-point insertion-loss curves. It is inspired by a spectral transposed
convolutional neural network (S-TCNN), but is not an exact reproduction of the
published architecture. The selected implementation contains 35,870 trainable
parameters and proceeds as follows:

\begin{enumerate}
    \item A 32-unit dense encoder maps the design vector to a latent state.
    \item A dense projection creates 800 values and reshapes them to 25
          positions with 32 channels.
    \item Three one-dimensional transposed convolutions double the sequence
          length from 25 to 50, 100 and 200 positions, with 32, 16 and 8
          channels respectively.
    \item Nine fixed frequency coordinates---normalised frequency plus four
          sine/cosine pairs---are concatenated at each output position.
    \item A one-dimensional refinement convolution and a linear six-channel
          convolution return the $200\times6$ insertion-loss tensor.
\end{enumerate}

\begin{figure*}[p]
    \centering
    \includegraphics[height=0.87\textheight,keepaspectratio]
      {appendix_d/curve_neural_topology.png}
    \caption{Topology of the selected whole-curve neural model. Transposed
    convolutions expand the latent sequence to 200 frequencies before the fixed
    frequency coordinates and final curve-refinement layers are applied.}
    \label{fig:appendix-d-curve-topology}
\end{figure*}

The selected curve loss combines point-wise MSE with error in the first
difference along frequency:
\begin{equation}
    \mathcal L_{\mathrm{curve}}
      =\operatorname{MSE}(\widehat{\mathbf L},\mathbf L)
      +\lambda_{\Delta}\operatorname{MSE}
       (\Delta_f\widehat{\mathbf L},\Delta_f\mathbf L).
    \label{eq:appendix-d-curve-loss}
\end{equation}
The training-derived weight $\lambda_{\Delta}=11.626038$ makes the derivative
term contribute 10\% of the point-wise MSE on the training reference. Validation
experiments compared absent, linear and Fourier frequency coordinates;
current-size and compact decoders; and point-wise and curve-aware losses. The
compact Fourier, curve-aware configuration was retained. Because interface,
scaling, architecture and loss changed together, its comparison with the
point-wise MLP is not a controlled ablation of frequency structure alone.

Training uses Adam with an initial learning rate of $10^{-3}$, batches of 64,
gradient clipping at 0.5 and a maximum of 100 epochs. Learning-rate reduction
and early stopping monitor inverse-scaled validation mean absolute error (MAE)
across all six curves.

\subsection{Reciprocal Complete-Matrix Neural Model}
The final model maps one design to the complete complex response at 200
frequencies. Its external data contract contains real and imaginary parts for
all $12^2=144$ entries, or 288 real values per frequency. Internally, it predicts
only the 78 unique upper-triangle complex entries, represented by 156 real
channels, and constructs the remaining entries by
\begin{equation}
    \widehat S_{i,j}(f_k)=\widehat S_{j,i}(f_k),\qquad i>j.
    \label{eq:appendix-d-reciprocity}
\end{equation}
This guarantees reciprocity exactly while avoiding duplicate learned outputs.
It does not enforce passivity or causality.

The network has 125,852 trainable parameters. The ten standardised design values
are repeated over the 200-point grid. At each frequency they are concatenated
with 41 fixed coordinates: normalised frequency, four Fourier sine/cosine pairs
and 32 Gaussian radial-basis features. A dense projection produces 128 features.
Three width-128 residual blocks then learn corrections while skip connections
preserve a direct route through each block. A final linear layer produces the
$200\times156$ scaled real/imaginary output.

\begin{figure*}[p]
    \centering
    \includegraphics[height=0.90\textheight,keepaspectratio]
      {appendix_d/full_smatrix_neural_topology.png}
    \caption{Topology of the selected frequency-conditioned full-matrix neural
    model. The graph ends at the 156 unique real-valued channels. Reciprocal
    mirroring and inverse scaling occur in the model wrapper after this Keras
    graph.}
    \label{fig:appendix-d-full-matrix-topology}
\end{figure*}

The selected loss combines normalised complex MSE with a weight-0.1 Huber loss
on off-diagonal log magnitude. The log-magnitude term reduces the dominance of
large responses and gives weak transmission/crosstalk entries a direct training
signal. The passivity weight is zero and there is no causality loss. Validation
complex normalised root-mean-square error (NRMSE) selects the restored
checkpoint. Adam training uses a $10^{-3}$ initial learning rate, batches of 64,
gradient clipping at 0.5 and the same maximum of 100 epochs with plateau-based
learning-rate reduction and early stopping.

The approved plan proposed a row-wise input $[\mathbf u,f]$ and $2m^2$ outputs
per design--frequency sample. The implemented model instead maps one design to
the whole fixed grid and predicts the reciprocal upper triangle. This is a
method change that retained the complete-matrix objective. Architectural
reciprocity is an additional constraint; it is not a substitute for the
planned passivity and causality enforcement.

\subsection{Alternatives and Engineering Decisions}
Table~\ref{tab:appendix-d-decisions} records the main alternatives that shaped
the implemented workflow. The decisions are supported by validation evidence
or by an explicit representation requirement, rather than by test-set ranking
alone.

\begin{table*}[t]
    \caption{Principal design decisions and retained boundaries.}
    \label{tab:appendix-d-decisions}
    \centering
    \scriptsize
    \begin{tabular}{@{}p{0.18\textwidth}p{0.25\textwidth}
                         p{0.25\textwidth}p{0.24\textwidth}@{}}
        \toprule
        Decision & Alternatives considered & Retained choice and reason
                 & Boundary \\
        \midrule
        Point-wise reference
          & scalar/vector Ridge, powers-only Ridge, forest
          & retain all as progressively different capacity checks
          & no model wins every error criterion \\
        Explicit powers for MLP
          & raw standardised inputs or fifth-degree powers
          & evaluate both with the same hidden layers
          & powers add no cross-feature products \\
        Curve coordinates
          & none, linear or Fourier
          & four Fourier pairs selected by validation
          & fixed grid; no frequency extrapolation test \\
        Curve decoder
          & current-size or compact; point-wise or derivative-aware loss
          & compact decoder with curve-aware loss selected by validation
          & several components change together \\
        Complex representation
          & magnitude/phase, complex-valued layers, or real/imaginary channels
          & real/imaginary channels preserve phase without custom complex layers
          & real-valued network does not itself ensure physical validity \\
        Matrix output
          & predict all 144 entries or only reciprocal upper triangle
          & predict 78 entries and mirror them to guarantee reciprocity
          & reciprocity only; no passivity/causality guarantee \\
        Physical loss
          & passivity penalty and additional weak-response terms
          & selected run uses log-magnitude weight 0.1 and passivity weight 0
          & enforcement/ablation remains incomplete \\
        \bottomrule
    \end{tabular}
\end{table*}

\subsection{Requirements-to-Verification Traceability}
Detailed numerical results are reported in Appendix~\ref{app:testing-results}.
Table~\ref{tab:appendix-d-traceability} connects the approved technical intent
to the implemented design and its verification evidence. It distinguishes
implemented checks from objectives that were not completed.

\begin{table*}[t]
    \caption{Technical requirements, implementation and verification.}
    \label{tab:appendix-d-traceability}
    \centering
    \scriptsize
    \begin{tabular}{@{}p{0.22\textwidth}p{0.30\textwidth}
                         p{0.38\textwidth}@{}}
        \toprule
        Technical intent & Implementation & Verification/evidence \\
        \midrule
        Reproducible parameter/Touchstone pipeline
          & identifier join, documented exclusions and design-level split
          & 7,030 aligned designs; saved split identifiers, preprocessors and
            run manifests \\
        Non-neural baseline stronger than a trivial mean
          & scalar and vector Ridge plus stronger mean-curve reference
          & global-mean part met; Scalar Ridge does not beat the
            frequency-dependent mean curve \\
        Broadband selected-path prediction
          & whole-design $10\rightarrow200\times6$ curve decoder
          & validation/test errors, paired design bootstrap and response plots
            in Appendix E \\
        Complete complex matrix prediction
          & $10\rightarrow200\times156$ unique channels, mirrored to
            $200\times12\times12$ complex matrices
          & same-task training-mean comparison and representative diagonal and
            off-diagonal responses \\
        Reciprocity
          & upper-triangle reconstruction
          & zero prediction reciprocity residual by construction \\
        Passivity and causality
          & post-training sampled passivity and finite-band causality diagnostics
          & no sampled passivity violations; causality is diagnostic only;
            planned enforcement and matched ablation not implemented \\
        Controlled parameter variation
          & no final controlled sweep
          & unmet; no sensitivity claim \\
        Reduced electromagnetic-simulation cost
          & model inference recorded without matched simulation timing
          & not established; no speed-up claim \\
        \bottomrule
    \end{tabular}
\end{table*}

\subsection{Selected Artifacts and Reproduction}
The selected neural diagrams were generated by Notebook 08 (NB08) from the
selected-model registry and saved Keras files; no model was retrained or
reselected. The machine-readable inventory is
\nolinkurl{project_portfolio/evidence/appendix_d_neural_model_inventory.csv}.
The four diagram files are stored under
\nolinkurl{project_portfolio/media/appendix_d/}. Non-neural selected estimators
are stored as joblib wrappers and neural estimators as Keras models together
with their preprocessing state, resolved configuration, training history,
metrics and environment record.

The selected full-matrix environment records CPython 3.10.20, TensorFlow 2.16.2,
Keras 3.12.2, NumPy 1.26.4, pandas 2.3.3 and scikit-learn 1.7.2 on Apple arm64.
NB08 additionally requires Graphviz and \texttt{pydot} to call
\texttt{keras.utils.plot\_model()}; both dependencies are declared in
\nolinkurl{sparam-surrogate/environment.yml}. Appendix E lists the selected run
identifiers and the independent metric-reproduction checks.

% Keep Appendix D's deferred full-width topology figures ahead of Appendix E.
\clearpage
